% ---------- Слайд 9: Формулы ----------
\begin{frame}[fragile,label=formulas]{Формулы и ссылки}
  \begin{columns}[T,onlytextwidth]
    \column{0.52\textwidth}
      \textbf{Исходный код:}
      \vspace{0.2em}
      \inputminted[fontsize=\scriptsize, frame=single, breaklines=true]{latex}{examples/formula.tex}

    \column{0.44\textwidth}
      \textbf{Результат:}
      \vspace{0.2em}
      
      \begin{equation} \label{eq1}
  E = mc^2
\end{equation}

\begin{equation} \label{eq2}
  \quad \int_{-\infty}^{\infty} e^{-x^2}\,dx = \sqrt{\pi}
\end{equation}

Ссылка на формулы по именам \ref{eq1} и \ref{eq2}.
  \end{columns}
  \note{Формулы из статьи переносятся в слайд без изменений.}
\end{frame}

% ---------- Слайд 10: Графика и таблицы ----------
\begin{frame}[fragile,label=graphics]{Графика и таблицы}
  \begin{columns}[T,onlytextwidth]
    % --- ЛЕВАЯ КОЛОНКА: КОД ---
    \column{0.52\textwidth}
      \textbf{Исходный код:}
      \vspace{0.2em}
      \inputminted[fontsize=\scriptsize, frame=single, breaklines=true]{latex}{examples/graphics.tex}

    % --- ПРАВАЯ КОЛОНКА: РЕЗУЛЬТАТ ---
    \column{0.44\textwidth}
      \textbf{Результат:}
      \vspace{0.2em}
      \begin{center}
  \includegraphics[width=0.6\linewidth] {example-image-a}
\end{center}

\begin{center}
  \begin{tabular}{lcc}
    \toprule
    Метод & Точность & Время \\
    \midrule
    A & 0.91 & 12\,с \\
    B & 0.94 & 17\,с \\
    \bottomrule
  \end{tabular}
\end{center}
  \end{columns}
  
  \vspace{0.2em}
  \footnotesize \emph{Рекомендация: графики лучше экспортировать в PDF (векторный формат).}
  \note{Пакет booktabs исключает вертикальные линии и использует
  горизонтальные разделители разной толщины.}
\end{frame}